% Bewertungspaket zum Gesamttest «Trigonometrische Funktionen» — Quelle des PDFs (python3 scripts/build-lp-pdf.py trigonometrische).
% Ganze Punkte je Teilschritt; Werte mit python3 nachgerechnet (05.10.2026, Folgefehler-Fälle eingeschlossen).
\documentclass[11pt]{article}
\usepackage{../lp-druck}
\renewcommand{\lpname}{Trigonometrische Funktionen}
\renewcommand{\lpbereich}{Schwerpunktfach}
\renewcommand{\lpteilgebiet}{3.5}
\renewcommand{\lpfuss}{Bewertungspaket}
\newcolumntype{P}{>{\centering\arraybackslash}p{10mm}}
\newenvironment{raster}{\par\smallskip\noindent\tabularx{\linewidth}{@{}>{\raggedright\arraybackslash}p{38mm}P X@{}}\toprule
  \small\textsc{Teilschritt} & \small\textsc{P} & \small\textsc{Musterlösung}\\\midrule}{\bottomrule\endtabularx\par}
\newcommand{\fehler}[1]{\par\smallskip{\small\textcolor{lprot}{\bfseries Typische Fehler:} #1}}
\newcommand{\E}{\,{\footnotesize\bfseries\color{lpgruen}(E)}}
\newcommand{\A}{\,{\footnotesize\bfseries\color{lpblau}(A)}}
\begin{document}
\lpkopf{Bewertungspaket zum Gesamttest}{}

\begin{lpachtung}{Erst lösen, dann öffnen}
Dieses Paket enthält die Musterlösung. Wer es vor dem Lösen liest, misst am Ende nur, wie gut das Abschreiben gelingt.
\end{lpachtung}

\section*{1 · So gehst du vor}
\begin{enumerate}[leftmargin=*,itemsep=2pt]
\item \textbf{Gesamttest lösen} — vollständig, mit Rechenweg, ohne dieses Paket; Teil A ohne, Teil B mit Taschenrechner.
\item \textbf{Fotografieren oder scannen}, gut lesbar, eine Seite pro Bild. \textbf{Kein Name} auf den Bildern.
  Handyfotos tragen oft unsichtbar Standort, Zeit und Gerät mit (Metadaten). Lade darum lieber einen Scan oder ein
  Bildschirmfoto des Fotos hoch, oder entferne vorher die Standortangaben.
\item \textbf{Bewerten lassen.} Ob du dafür eine KI nutzt und welche, entscheidest du selbst und in eigener Verantwortung —
  je nachdem, was dir aufgrund deines Alters und deiner persönlichen Situation erlaubt ist (Altersgrenzen und
  Nutzungsbedingungen des Anbieters, allenfalls Einverständnis deiner Eltern). Lade nur deine Lösung und dieses PDF hoch,
  keine persönlichen Daten, und füge den Auftrag aus Abschnitt~2 ein. Ohne KI kannst du dich mit Abschnitt~3 selbst bewerten.
\item \textbf{Rückmeldung prüfen.} Stimmt, was die KI aus deiner Handschrift gelesen hat? Eine KI kann sich irren —
  im Zweifel gilt das Raster in Abschnitt~3.
\item \textbf{Gezielt wiederholen} nach Abschnitt~4.
\end{enumerate}

\needspace{14\baselineskip}
\section*{2 · Auftrag an die KI}
\begin{tcolorbox}[colback=lplinie!25,colframe=lplinie,boxrule=0.5pt,arc=1mm,fontupper=\small]
Du bist Korrektorin oder Korrektor für Mathematik an einer Schweizer Berufsmaturitätsschule. Im Anhang findest du (1)~das Bewertungspaket zum Gesamttest «Trigonometrische Funktionen» mit Musterlösung und Punkteraster und (2)~Fotos meiner handschriftlichen Lösung.\par\medskip
Geh so vor:\par
1. Schreib zu jeder Aufgabe G1 bis G8 zuerst ab, was du in meiner Lösung liest — Rechenweg und Ergebnis. Was du nicht sicher lesen kannst, markierst du mit [unsicher]. Nicht raten.\par
2. Vergib die Punkte genau nach dem Raster im Bewertungspaket (Abschnitt 3), ganze Punkte je Zeile. Ziehe einen Fehler nur einmal ab: Rechne ich mit einem falschen Zwischenergebnis richtig weiter, gibt es die Folgepunkte. Ausnahme: Zeilen mit \textbf{(E)} gibt es nur für das richtige Ergebnis, nie als Folgepunkt. Die Zeilen «Typische Fehler» sagen, welche Rasterzeile bei diesem Fehler 0 Punkte gibt — sie sind kein zusätzlicher Abzug.\par
3. Andere richtige Lösungswege zählen voll. Gleichwertige Schreibweisen zählen voll: Dezimalpunkt oder Dezimalkomma, Bruch oder Dezimalzahl ($\tfrac18 = 0.125$), $y = \ldots$ oder $f(x) = \ldots$, Bogenmass exakt oder auf drei Dezimalen ($\tfrac{5\pi}{6}$ oder $2.618$), $k \in \mathbb{Z}$ oder «$k$ ganz» — ausser wo das Raster eine bestimmte Form verlangt. Drei Stufen: Zeilen mit \textbf{(A)} verlangen nur Ablesen, diese Punkte gibt es auch ohne Rechenweg. Zeilen mit \textbf{(E)} sind Ergebnispunkte: Das Ergebnis muss stimmen \emph{und} aus einem erkennbaren Weg hervorgehen, ausser die Zeile sagt ausdrücklich etwas anderes. Alle übrigen Zeilen bewerten den Weg selbst.\par
4. Begründe jeden Abzug in einem Satz und nenne die Fehlerart (zum Beispiel «Periode mal statt durch b» oder «nur die Lösung des Rechners»). Schreib die Musterlösung nicht ab — zeig mir, wo mein Fehler liegt.\par
5. Gib zum Schluss eine Tabelle «Aufgabe | Punkte | Begründung», die Gesamtpunktzahl (von 25) und — nach der Selbsteinschätzung in Abschnitt 4 — welches Kapitel des Leitprogramms ich wiederholen soll. Keine Note.\par
6. Fehlt ein Foto oder ist eine Aufgabe unlesbar, sag es, statt sie zu bewerten.
\end{tcolorbox}

\needspace{16\baselineskip}
\section*{3 · Musterlösung und Punkteraster}
\textbf{Allgemein:} Ganze Punkte je Zeile. Ein Fehler wird nur einmal abgezogen (Folgefehler). Gleichwertige Wege und
Schreibweisen zählen voll. In der Spalte P:
\textbf{(A)}~= nur ablesen, zählt auch ohne Rechenweg · \textbf{(E)}~= Ergebnispunkt, Ergebnis richtig
\emph{und} Weg erkennbar (ausser die Zeile sagt ausdrücklich etwas anderes) · ohne Zeichen = die Zeile bewertet den Weg.
Teil A ist ohne, Teil B mit Taschenrechner gestellt.
«Typische Fehler» nennt, welche Zeile 0 Punkte gibt, und die Punkte, die bleiben.
\textbf{Teilrichtige Zeilen:} Verlangt eine Zeile mehrere Angaben und stimmt nur ein Teil, gibt die Zeile 0 Punkte —
ausser die Zeile nennt ausdrücklich eine Aufteilung.

\begin{lpaufgabe}{G1}{Sinus und Cosinus skizzieren}{3 P}
\begin{raster}
(a) Sinuskurve & 1 & durch $(-\pi \mid 0)$, $\left(-\tfrac{\pi}{2} \mid -1\right)$, $(0 \mid 0)$, $\left(\tfrac{\pi}{2} \mid 1\right)$, $(\pi \mid 0)$, $\left(\tfrac{3\pi}{2} \mid -1\right)$, $(2\pi \mid 0)$; weiche Bögen, keine Spitzen. Toleranz: ein halbes Kästchen.\\
(a) Cosinuskurve & 1 & durch $(-\pi \mid -1)$, $\left(-\tfrac{\pi}{2} \mid 0\right)$, $(0 \mid 1)$, $\left(\tfrac{\pi}{2} \mid 0\right)$, $(\pi \mid -1)$, $\left(\tfrac{3\pi}{2} \mid 0\right)$, $(2\pi \mid 1)$.\\
(b) Versatz & 1\E & die Sinuskurve: $\cos\left(x - \tfrac{\pi}{2}\right) = \sin x$ — der Hochpunkt wandert von $0$ nach $\tfrac{\pi}{2}$, wo die Sinuskurve ihren hat. Gleichwertig: «$\sin x = \cos\left(x - \tfrac{\pi}{2}\right)$». Kurve und Gleichung nötig.\\
\end{raster}
\fehler{Sinus und Cosinus vertauscht: beide Skizzen-Zeilen 0 $\to$ höchstens 1 von 3\,P.
Nur $0 \le x \le 2\pi$ gezeichnet: die betroffene Skizzen-Zeile 0.
(b) «$-\sin x$» oder «$\cos x + \tfrac{\pi}{2}$» (nach oben statt nach rechts verschoben): Zeile (b) 0.}
\end{lpaufgabe}

\begin{lpaufgabe}{G2}{Werte ohne Rechner}{3 P}
\begin{raster}
(a) Grad & 1\E & $\tfrac{5\pi}{3} = \tfrac53 \cdot 180^\circ = 300^\circ$\\
(b) Sinus, Cosinus & 1\E & $\sin \tfrac{11\pi}{6} = -\tfrac12$ (vierter Quadrant), $\cos \tfrac{4\pi}{3} = -\tfrac12$ (dritter Quadrant). Beide nötig.\\
(c) Tangens & 1\E & $\tan \tfrac{3\pi}{2}$ ist nicht definiert ($\cos \tfrac{3\pi}{2} = 0$, Pol); $\tan\left(-\tfrac{3\pi}{4}\right) = 1$ (Sinus und Cosinus beide negativ, gleich gross). Beide nötig.\\
\end{raster}
\fehler{(b) $\tfrac12$ statt $-\tfrac12$ (Vorzeichen des Quadranten übersehen): Zeile (b) 0.
(c) $\tan \tfrac{3\pi}{2} = 0$ (Pol mit Nullstelle verwechselt) oder $\tan\left(-\tfrac{3\pi}{4}\right) = -1$: Zeile (c) 0.
Als (E)-Weg genügt der Quadrant oder ein Hinweis auf den Einheitskreis.}
\end{lpaufgabe}

\begin{lpaufgabe}{G3}{Symmetrien}{3 P}
\begin{raster}
(a) & 1\E & $\sin(-1.1) = -\sin 1.1 \approx -0.891$ — Sinuskurve punktsymmetrisch zum Ursprung.\\
(b) & 1\E & $\cos(-1.1) = \cos 1.1 \approx 0.454$ — Cosinuskurve achsensymmetrisch zur $y$-Achse.\\
(c) & 1\E & $\sin(\pi - 1.1) = \sin 1.1 \approx 0.891$ — Sinuskurve symmetrisch zur Geraden $x = \tfrac{\pi}{2}$.\\
\end{raster}
\fehler{Vorzeichen in (a) vergessen ($0.891$): Zeile (a) 0. (b) $-0.454$: Zeile (b) 0. (c) $\cos 1.1 \approx 0.454$ oder $-0.891$: Zeile (c) 0.
Richtige Zahl ohne jede Begründung: die Zeile 0 — (E) verlangt einen erkennbaren Weg, hier die Symmetrie.}
\end{lpaufgabe}

\begin{lpaufgabe}{G4}{Die Tangenskurve}{3 P}
\begin{raster}
Polgeraden & 1\A & senkrecht bei $x = \tfrac{\pi}{2}$ und $x = \tfrac{3\pi}{2}$, beide eingezeichnet (gestrichelt oder durchgezogen).\\
Kurve & 1 & drei Äste, steigend; durch die Nullstellen $(0 \mid 0)$, $(\pi \mid 0)$, $(2\pi \mid 0)$ und z.\,B. $\left(\tfrac{\pi}{4} \mid 1\right)$, $\left(\tfrac{3\pi}{4} \mid -1\right)$; nähert sich den Polgeraden, ohne sie zu berühren.\\
$p$ und $W$ & 1\A & $p = \pi$, $W = \mathbb{R}$. Beides nötig.\\
\end{raster}
\fehler{Pole bei $0$, $\pi$, $2\pi$ (mit den Nullstellen vertauscht): Polgeraden-Zeile und Kurven-Zeile 0 $\to$ höchstens 1 von 3\,P.
Äste fallend gezeichnet oder Kurve über die Polgerade hinweg verbunden: Kurven-Zeile 0.
$p = 2\pi$ oder $W = [-1;\, 1]$: letzte Zeile 0.}
\end{lpaufgabe}

\begin{lpaufgabe}{G5}{Schrittweise skizzieren}{3 P}
\begin{raster}
$p$ und $W$ & 1\E & $p = \tfrac{2\pi}{3}$; $W = [0.5 - 2;\, 0.5 + 2] = [-1.5;\, 2.5]$. Beides nötig.\\
Hochpunkt & 1\E & $3\left(x - \tfrac{\pi}{6}\right) = \tfrac{\pi}{2} \Rightarrow x = \tfrac{\pi}{6} + \tfrac{\pi}{6} = \tfrac{\pi}{3}$: $H\left(\tfrac{\pi}{3} \mid 2.5\right)$\\
Skizze & 1 & Mittellinie $y = 0.5$; steigend durch $\left(\tfrac{\pi}{6} \mid 0.5\right)$, Hochpunkte bei $\tfrac{\pi}{3}$ und $\pi$ ($y = 2.5$), Tiefpunkte bei $0$, $\tfrac{2\pi}{3}$ und $\tfrac{4\pi}{3}$ ($y = -1.5$). Toleranz: ein halbes Kästchen.\\
\end{raster}
\fehler{$p = 6\pi$ ($2\pi \cdot 3$): Zeile $p$/$W$ 0; die Skizze mit falscher Periode ebenfalls 0 $\to$ höchstens 1 von 3\,P.
Nach links verschoben gerechnet, $3\left(x + \tfrac{\pi}{6}\right) = \tfrac{\pi}{2} \Rightarrow x = 0$: Hochpunkt-Zeile 0.
Hochpunkt bei $\tfrac{\pi}{6} + \tfrac{\pi}{2} = \tfrac{2\pi}{3}$ (Stauchung durch $b = 3$ vergessen): Hochpunkt-Zeile 0 — erst $3\left(x - \tfrac{\pi}{6}\right) = \tfrac{\pi}{2}$ auflösen.}
\end{lpaufgabe}

\begin{lpaufgabe}{G6}{Gleichung aus dem Graphen}{3 P}
\begin{raster}
$v$ & 1\A & höchster Wert $3.5$, tiefster $-1.5$ (beide an der Achse beschriftet): $v = \tfrac{3.5 + (-1.5)}{2} = 1$\\
$a$ & 1\A & $a = 3.5 - 1 = 2.5$\\
$b$ & 1\E & Periode $p = \pi$ (von Hochpunkt $\tfrac{\pi}{4}$ zu Hochpunkt $\tfrac{5\pi}{4}$): $b = \tfrac{2\pi}{\pi} = 2$. Also $y = 2.5\sin(2x) + 1$.\\
\end{raster}
\fehler{$b = 0.5$ (Periode richtig abgelesen, aber $b$ und $\tfrac{1}{b}$ verwechselt): Zeile $b$ 0 $\to$ 2 von 3\,P.
$a = 5$ (ganze Spannweite statt halbe): Zeile $a$ 0 $\to$ 2 von 3\,P.
$v = 0$ angenommen und darum $a = 3.5$: Zeilen $v$ und $a$ je 0 (beide sind Ablesezeilen und gelten nur mit dem richtigen Wert) $\to$ höchstens 1 von 3\,P.}
\end{lpaufgabe}

\begin{lpaufgabe}{G7}{Alle Lösungen}{4 P}
\begin{raster}
(a) $x_1$ & 1\E & $x_1 = \sin^{-1}(0.35) \approx 0.358$\\
(a) $x_2$ & 1 & Regel $x_2 = \pi - x_1$ angewandt: $\approx 2.784$. Der Punkt gilt auch mit einem falschen $x_1$, wenn die Regel stimmt.\\
(b) $x_1$ & 1\E & $x_1 = \cos^{-1}(0.35) \approx 1.213$\\
(b) $x_2$ & 1 & Regel $x_2 = 2\pi - x_1$ angewandt: $\approx 5.070$. Der Punkt gilt auch mit einem falschen $x_1$, wenn die Regel stimmt.\\
\end{raster}
\fehler{Rechner im Gradmodus ($20.487$ bzw. $69.513$): beide $x_1$-Zeilen 0. Wird die Regel in Grad richtig angewandt ($180^\circ - x_1$ bzw. $360^\circ - x_1$), gibt es die $x_2$-Zeilen $\to$ 2 von 4\,P; sonst 0 von 4\,P.
(a) $x_2 = 2\pi - x_1 \approx 5.925$ (Regel des Cosinus): Zeile (a) $x_2$ 0.
(b) $x_2 = \pi - x_1 \approx 1.928$ (Regel des Sinus): Zeile (b) $x_2$ 0.
Nur die Rechnerwerte, keine zweiten Lösungen $\to$ 2 von 4\,P. Abweichung in der dritten Dezimale durch Runden zählt voll.}
\end{lpaufgabe}

\begin{lpaufgabe}{G8}{Ebbe und Flut}{3 P}
\begin{raster}
(a) Extremwerte & 1\A & höchstens $3 + 2 = 5$ m, mindestens $3 - 2 = 1$ m. Beides nötig.\\
(b) Periode & 1\E & $p = \tfrac{2\pi}{\pi/6} = 12$ Stunden\\
(c) Zeiten & 1\E & $2\sin\left(\tfrac{\pi}{6}t\right) + 3 = 4 \Rightarrow \sin\left(\tfrac{\pi}{6}t\right) = \tfrac12 \Rightarrow \tfrac{\pi}{6}t = \tfrac{\pi}{6}$ oder $\tfrac{5\pi}{6}$, also $t = 1$ und $t = 5$ (um $1$ Uhr und um $5$ Uhr). Beide nötig.\\
\end{raster}
\fehler{(b) $p = \tfrac{\pi^2}{3}$ ($2\pi$ mal statt durch $b$): Zeile (b) 0.
(c) Nur $t = 1$ (zweite Lösung über die Symmetrie vergessen): Zeile (c) 0 $\to$ 2 von 3\,P.
(c) Mit dem Rechner $t = \tfrac{6}{\pi}\sin^{-1}(0.5) = 1$ und $t = 6 - 1 = 5$: gleichwertig, zählt voll.}
\end{lpaufgabe}

\needspace{14\baselineskip}
\section*{4 · Selbsteinschätzung}
\begin{tabularx}{\linewidth}{@{}l X@{}}\toprule
22 – 25 P & Die geprüften Teile sitzen. Wo du Punkte verloren hast: das Kapitel dieser Aufgabe nochmals (Zuordnung unten).\\
17 – 21 P & Den schwächsten Teil nochmals: Simulation und Übungen des Kapitels, in dem du die meisten Punkte verloren hast.\\
11 – 16 P & Zurück zu den Kapiteln aller Aufgaben, in denen du Punkte verloren hast.\\
0 – 10 P & Zurück zu Kapitel 1 und von dort der Reihe nach weiter.\\\midrule
\multicolumn{2}{@{}p{\linewidth}@{}}{\small\textbf{Aufgabe $\to$ Kapitel:} G1 $\to$ Kapitel 1, 2; G2 $\to$ Kapitel 1, 3; G3 $\to$ Kapitel 2; G4 $\to$ Kapitel 3; G5, G6 $\to$ Kapitel 4; G7 $\to$ Kapitel 5; G8 $\to$ Kapitel 4, 5}\\\bottomrule
\end{tabularx}
\end{document}
